\subsection*{Algorithm Track Benchmarks}

    The v1.0 iteration of the NeuroBench algorithm track includes four benchmarks for neuromorphic computing research. The benchmarks were chosen by the NeuroBench community to capture key ongoing challenges for neuromorphic algorithm design. The list of tasks highlights features which are relevant to neuromorphic research interests: few-shot continual learning, object detection utilizing the high dynamic range and temporal resolution of event cameras, sensorimotor decoding based on cortical signals, and low-dimensional predictive modeling useful for prototyping resource-constrained networks that are suitable for small mixed-signal systems.
    % The benchmark tasks and baselines were produced by working groups comprised of various industry and academic researchers within the NeuroBench community. 
    Benchmark tasks are listed below and summarized in Table \ref{tab:alg_benchmarks}. Detailed specifications of benchmark tasks are provided in the Methods section. 
    
    % \begin{table}[t]
%         \centering
%         { \renewcommand{\arraystretch}{1.2}
%         \begin{tabular}{|c|c|c|p{0.23\linewidth}|}
%         \hline
%              Task & Dataset & Correctness metric(s) & Task description  \\\hline
%              Keyword FSCIL & MSWC~\cite{mazumder21mswc} & Accuracy & Few-shot, continual learning of keyword classes. \\\hline
%              Event Camera Object Detection & Prophesee 1MP Automotive~\cite{Perot2020}  & COCO mAP & Detecting automotive objects from event-sensor video.  \\\hline
%              Motor Prediction & Primate Reaching~\cite{makin-dataset} & R$^2$ & Predicting fingertip velocity from cortical data.  \\\hline
%              Chaotic Function Prediction & Mackey-Glass time series~\cite{mackey1977oscillation} & sMAPE & Autoregressive modeling of chaotic functions.  \\\hline
              
%         \end{tabular}
%         }
%         \caption{NeuroBench algorithm track v1.0 benchmarks.}
%         \label{tab:alg_benchmarks}
%     \end{table}

\begin{table}[t]
    \centering
    \renewcommand{\arraystretch}{1.3}
    \begin{tabular}{lllm{0.23\linewidth}}
        \thickhline
        Task & Dataset & Correctness metric & Task description  \\\hline
        Keyword FSCIL & MSWC~\cite{mazumder21mswc} & Accuracy & Few-shot, continual learning of keyword classes. \\\hline
        Event Camera Object Detection & Prophesee 1MP Automotive~\cite{Perot2020}  & COCO mAP & Detecting automotive objects from event camera video.  \\\hline
        NHP Motor Prediction & Primate Reaching~\cite{makin-dataset} & R$^2$ & Predicting fingertip velocity from cortical recordings.  \\\hline
        Chaotic Function Prediction & Mackey-Glass time series~\cite{mackey1977oscillation} & sMAPE & Autoregressive modeling of chaotic functions.  \\\thickhline
    \end{tabular}
    \caption{NeuroBench algorithm track v1.0 benchmarks.}
    \label{tab:alg_benchmarks}
\end{table}


    
    \begin{itemize}
        \item \textbf{Keyword Few-Shot Class-Incremental Learning (FSCIL)} --
        Learning new tasks from a small amount of experiences while retaining knowledge of prior tasks is a hallmark of biological intelligence and a long-standing goal of general AI~\cite{kudithipudi22biological}. It is especially a key challenge to endow edge devices with the ability to adapt to their environments and users. This benchmark thus evaluates the capacity of a model to successively incorporate new keywords over multiple sessions (class-incremental), with only a handful of samples from the new classes to train with (few-shot). The FSCIL task is a recently established benchmark in the computer vision domain~\cite{tao2020fscil}, but it has not yet been adapted to other data modalities. Aligning with a neuromorphic interest in temporal data modalities, this benchmark introduces a FSCIL task with streaming audio data using the large Multilingual Spoken Word Corpus (MSWC)~\cite{mazumder21mswc} keyword classification dataset. The task is designed to be approached in two phases: pre-training and incremental learning. First, for pre-training, a set of 100 words spanning 5 base languages (English, German, Catalan, French, Kinyarwanda) with 500 training samples each are made available to train an initial model. Next, for incremental learning, the model undergoes 10 successive sessions to learn words from 10 new languages (Persian, Spanish, Russian, Welsh, Italian, Basque, Polish, Esparanto, Portuguese, Dutch) in a few-shot learning scenario. Each incremental session adds 10 words of the corresponding session language with only 5 training samples available per word. After each session, the model is tested in classification accuracy on all prior learned classes, including the 100 base pre-training classes and the few-shot-learned classes,
        %, with 100 test samples per class
        therefore evaluating the FSCIL solution on its ability to learn new classes while retaining knowledge about the previously learned ones. Each session learns a new language, for a total knowledge base of 200 keywords by the end of the benchmark.

        
        \item \textbf{Event Camera Object Detection} --
        Object detection is a widely-used computer vision task with applications in robotics, autonomous driving, and surveillance. Such scenarios at the edge may require high energy efficiency and real-time performance, which can be achieved via event-based vision sensors~\cite{Gallego2022}. The event camera object detection benchmark uses the Prophesee 1 Megapixel automotive detection dataset~\cite{Perot2020}, a large labeled object detection dataset with over 15 hours of event camera video from the front of a car driving in various scenarios. Predetermined training, validation, and testing splits include \SI{11.2}{\hour}, \SI{2.2}{\hour}, and \SI{2.2}{\hour} of recording, respectively. Pedestrian, two-wheeler, and car object classes are used in evaluation, and correctness is measured using COCO mean average precision (mAP)~\cite{LinCOCO_2014}.
        
        \item \textbf{Non-human Primate (NHP) Motor Prediction} --
        Studying models which can accurately replicate features of biological computation presents opportunities in understanding sensorimotor behavior and developing closed-loop methods for future robotic agents. It also is foundational to the development of wearable or implantable neuro-prosthetic devices that can accurately generate motor activity from neural or muscle signals. This benchmark utilizes a dataset consisting of multi-channel recordings from the sensorimotor cortex of two non-human primates (NHP Indy and NHP Loco) during reaching movements, along with corresponding fingertip motion of the reach~\cite{makin-dataset}. Six total sessions are included from the dataset, for a total of 8712 seconds of data. The task is to train a model to predict the two-dimensional components of finger velocity using recent neural data. The sessions are treated independently (i.e., models are trained separately for each session), and the data is split to allow the first 75\% for training and validation and the last 25\% for evaluation. Correctness of the predictions is evaluated by the coefficient of determination ($R^{2}$) score against the true finger velocity targets, averaged over all six sessions.
        
        \item \textbf{Chaotic Function Prediction} --
        The real-world data benchmarks presented thus far are high-dimensional and can require large networks to achieve high accuracy, raising challenges for solution types with limited I/O support and network capacity, such as mixed-signal edge prototype solutions. To address this, we include a synthetic benchmark based on prediction of one-dimensional Mackey-Glass time series~\cite{mackey1977oscillation}, which can be effectively tackled by smaller networks. Mackey-Glass has been widely adopted as a benchmark for evaluating temporal predictors, including neuromorphic models~\cite{jaeger2004harnessing, mukhopadhyay2020learning, chilkuri2021parallelizing}. The task involves prediction of the next timestep value $f(t+\Delta t)$ given the current timestep value $f(t)$. 
        The model is trained and validated using the first half of the time series, during which the ground truth state $f(t)$ are supplied to the model to predict the next timestep $f'(t+\Delta t)$.
        During the evaluation, the model uses its prior prediction $f'(t)$ to generate each next value $f'(t+\Delta t)$, autoregressively forecasting the second half of the time series.
        Correctness is measured using symmetric mean absolute percentage error (sMAPE) of the generated time series against the target time series, a standard metric in forecasting~\cite{MAKRIDAKIS202054}.
        The benchmark includes a set of 14 Mackey-Glass time series, which vary by the equation parameter $\tau$, the delay constant. Lyapunov time $(L)$, the expected predictability timescale for chaos~\cite{gilpin2023largescale}, is used as the time unit for each time series. The total length of each series is 20 Lyapunov times, and 75 points are sampled per Lyapunov time ($\Delta t$ = $L/75$).

    \end{itemize}
